\section{\emph{E-Government}}
\label{sec:e-government}

\emph{Governance} merupakan proses lembaga-lembaga publik dalam mengatasi masalah-masalah publik, mengelola sumber daya publik, dan menjamin realisasi hak asasi manusia secara inklusif. Secara umum \emph{E-Government} dapat dikatakan sebagai suatu aplikasi berbasis komputer dan internet yang digunakan untuk meningkatkan hubungan dan layanan pemerintah kepada warga masyarakatnya. \emph{E-Government} memanfaatan teknologi informasi untuk meningkatkan hubungan dengan pihak-pihak dalam aspek \emph{good governance} (masyarakat dan lembaga bisnis) dengan tujuan meningkatkan kualitas pelayanan yang efektif dan efisien \cite{muliawaty2020egov}.

Proyek MonPD Reborn merupakan bagian dari inisiatif penerapan \emph{E-Government} di lingkungan Pemerintah Kota Surabaya, khususnya pada kategori G2G (\emph{Government-to-Government}). Aplikasi ini tidak berinteraksi langsung dengan masyarakat, melainkan menjadi alat kerja internal bagi para pegawai Bapenda. Dengan menyediakan platform terpusat untuk \emph{monitoring} data pajak, MonPD Reborn mendukung terciptanya pemerintahan yang lebih efisien, responsif, dan akuntabel, sejalan dengan visi Surabaya sebagai \emph{Smart City}.
